\documentclass{article}
\usepackage[backend=biber,natbib=true,style=alphabetic,maxbibnames=50]{biblatex}
\addbibresource{/home/nqbh/reference/bib.bib}
\usepackage[utf8]{vietnam}
\usepackage{tocloft}
\renewcommand{\cftsecleader}{\cftdotfill{\cftdotsep}}
\usepackage[colorlinks=true,linkcolor=blue,urlcolor=red,citecolor=magenta]{hyperref}
\usepackage{amsmath,amssymb,amsthm,enumitem,float,graphicx,mathtools,tikz}
\usetikzlibrary{angles,calc,intersections,matrix,patterns,quotes,shadings}
\allowdisplaybreaks
\newtheorem{assumption}{Assumption}
\newtheorem{baitoan}{Bài toán}
\newtheorem{cauhoi}{Câu hỏi}
\newtheorem{conjecture}{Conjecture}
\newtheorem{corollary}{Corollary}
\newtheorem{dangtoan}{Dạng toán}
\newtheorem{definition}{Definition}
\newtheorem{dinhly}{Định lý}
\newtheorem{dinhnghia}{Định nghĩa}
\newtheorem{example}{Example}
\newtheorem{ghichu}{Ghi chú}
\newtheorem{hequa}{Hệ quả}
\newtheorem{hypothesis}{Hypothesis}
\newtheorem{lemma}{Lemma}
\newtheorem{luuy}{Lưu ý}
\newtheorem{nhanxet}{Nhận xét}
\newtheorem{notation}{Notation}
\newtheorem{note}{Note}
\newtheorem{principle}{Principle}
\newtheorem{problem}{Problem}
\newtheorem{proposition}{Proposition}
\newtheorem{question}{Question}
\newtheorem{remark}{Remark}
\newtheorem{theorem}{Theorem}
\newtheorem{vidu}{Ví dụ}
\usepackage[margin=2cm]{geometry}
\def\labelitemii{$\circ$}
\DeclareRobustCommand{\divby}{%
	\mathrel{\vbox{\baselineskip.65ex\lineskiplimit0pt\hbox{.}\hbox{.}\hbox{.}}}%
}
\setlist[itemize]{leftmargin=*}
\setlist[enumerate]{leftmargin=*}

\title{Programming Problems in Mathematical Analysis -- Bài Tập Lập Trình Giải Tích}
\author{Nguyễn Quản Bá Hồng\footnote{A scientist- {\it\&} creative artist wannabe, a mathematics {\it\&} computer science lecturer of Department of Artificial Intelligence {\it\&} Data Science (AIDS), School of Technology (SOT), UMT Trường Đại học Quản lý {\it\&} Công nghệ TP.HCM, Hồ Chí Minh City, Việt Nam.\\E-mail: {\sf nguyenquanbahong@gmail.com} {\it\&} {\sf hong.nguyenquanba@umt.edu.vn}. Website: \url{https://nqbh.github.io/}. GitHub: \url{https://github.com/NQBH}.}}
\date{\today}

\begin{document}
\maketitle
\begin{abstract}
	This text is a part of the series {\it Some Topics in Advanced STEM \& Beyond}:
	
	{\sc url}: \url{https://nqbh.github.io/advanced_STEM/}.
	
	Latest version:
	\begin{itemize}
		\item {\it Programming Problems in Mathematical Analysis -- Bài Tập Lập Trình Giải Tích}.
		
		PDF: {\sc url}: \url{.pdf}.
		
		\TeX: {\sc url}: \url{.tex}.
		\item {\it }.
		
		PDF: {\sc url}: \url{.pdf}.
		
		\TeX: {\sc url}: \url{.tex}.
	\end{itemize}
\end{abstract}
\tableofcontents

%------------------------------------------------------------------------------%

\section{Sequences}

\begin{baitoan}
	Viết chương trình C++ để kiểm tra 1 dãy số thực $\{a_n\}_{n=0}^\infty\subset\mathbb{R}$ thỏa \cite[Def. 6.1.3]{Tao_analysis_1}.
\end{baitoan}

\begin{baitoan}
	Cho dãy số thực $\{a_n\}_{n=0}^\infty\subset\mathbb{R}$ được xác định bởi hàm
	\begin{verbatim}
double f(int  n) {
    return <formula of f(n) in C++>
}
	\end{verbatim}
	Tính $N_\varepsilon$ với $\varepsilon\in\mathbb{R}_{>0}$ bất kỳ bằng cách viết hàm
	\begin{verbatim}
int find_N_epsilon(vector<double> a, double epsilon) {
    // C++ code to find N_\epsilon
}
	\end{verbatim}
\end{baitoan}

\begin{baitoan}
	Viết chương trình C++ để kiểm tra 1 dãy số thực $\{a_n\}_{n=0}^\infty\subset\mathbb{R}$ thỏa \cite[Def. 6.1.5]{Tao_analysis_1}.
\end{baitoan}

\begin{baitoan}
	Viết chương trình C++ mô phỏng lại cách tìm supremum \& infimum của 1 dãy số thực $\{a_n\}_{n=1}^\infty$ được cài đặt bởi hàm sau
	\begin{verbatim}
double a(int m, int n) {
    return a_n := f(n)
}
	\end{verbatim}
	với $f:\mathbb{N}^\star\to\mathbb{R}$ là 1 hàm bất kỳ được viết riêng.
\end{baitoan}

\begin{proof}[Giải]
	
\end{proof}

%------------------------------------------------------------------------------%

\section{Miscellaneous}

%------------------------------------------------------------------------------%

\printbibliography[heading=bibintoc]
	
\end{document}